% Part: first-order-logic
% Chapter: axiomatic-deduction
% Section: proof-theoretic-notions

\documentclass[../../../include/open-logic-section]{subfiles}

\begin{document}

\iftag{FOL}
      {\olfileid{fol}{axd}{ptn}}
      {\olfileid{pl}{axd}{ptn}}

\olsection{ثبوتی-نظریاتی تصورات (Proof-Theoretic Notions)}

\begin{explain}
جیویں اسیں کئی اہم معنوی تصورات
(\iftag{FOL}{منطقی اعتبار}{ہمیشہ صادق ہونا}، منطقی نتیجہ، پورا ہون
جوگتا) دی تعریف دتی اے، ہُن اسیں اوہناں دے مطابق
\emph{ثبوتی-نظریاتی تصورات} دی تعریف دیندے آں۔ ایہناں دی تعریف
!!{sentence}s دے !!{structure}s وچ پورا اترن ول رجوع کر کے نہیں،
بلکہ کجھ فارمولیاں دے !!{derivability} یا !!{nonderivability} ول
رجوع کر کے دتی جاندی اے۔ ایہناں دوہاں طرح دے تصورات دا اک دوجے
نال میل کھانا اک اہم دریافت سی۔ ایس میل دا بیان \emph{صحت مندی}
تے \emph{تمامیت دے قضیاں} وچ اے۔
\end{explain}

\begin{defn}[!!^{derivability}]
!!^a{formula}~$!A$، $\Gamma$ توں \emph{!!{derivable}} اے، جیہنوں
$\Gamma \Proves !A$ لکھدے آں، جے~$\Gamma$ توں !!a{derivation}
موجود ہووے جیہڑا $!A$ اُتے ختم ہووے۔
\end{defn}

\begin{defn}[قضیے]
!!^a{formula}~$!A$ اک \emph{قضیہ} اے جے خالی سیٹ توں $!A$ دا
!!a{derivation} موجود ہووے۔ اسیں $\Proves !A$ لکھدے آں جے $!A$
قضیہ اے، تے $\Proves/ !A$ لکھدے آں جے نہیں۔
\end{defn}

\begin{defn}[بے تضادی]
!!{formula}s دا سیٹ~$\Gamma$ \emph{بے تضاد} اے اُدوں تے صرف اُدوں
جدوں $\Gamma\Proves/ \lfalse$؛ نہیں تاں ایہ \emph{متضاد} اے۔
\end{defn}

\begin{prop}[انعکاسیت]
\ollabel{prop:reflexivity}
جے $!A \in \Gamma$، تاں $\Gamma \Proves !A$۔
\end{prop}

\begin{proof}
  !!{formula}~$!A$ اپنے آپ وچ~$!A$ دا~$\Gamma$ توں
  !!a{derivation} اے۔
\end{proof}

\begin{prop}[یک رُخیت]
\ollabel{prop:monotonicity}
جے $\Gamma \subseteq \Delta$ تے $\Gamma \Proves !A$، تاں $\Delta
\Proves !A$۔
\end{prop}

\begin{proof}
$!A$ دا~$\Gamma$ توں کوئی وی !!{derivation}، $!A$ دا~$\Delta$
توں وی !!a{derivation} اے۔
\end{proof}

\begin{prop}[تعدیت]
\ollabel{prop:transitivity}
جے $\Gamma \Proves !A$ تے $\{!A\} \cup \Delta \Proves
!B$، تاں $\Gamma \cup \Delta \Proves !B$۔
\end{prop}

\begin{proof}
  فرض کرو $\{!A\} \cup \Delta \Proves !B$۔ تاں اک
  !!a{derivation} $!B_1$، \dots، $!B_l = !B$،~$\{!A\} \cup
  \Delta$ توں موجود اے۔ اوس !!{derivation} دے کجھ قدم ایسے
  قاعدے دی وجہ توں درست ہون گے جیہڑا پہلے آ چکی سطر~$!B_i = !A$
  دا حوالہ دیندا اے۔ مفروضے مطابق $!A$ دا~$\Gamma$ توں
  !!a{derivation} موجود اے، یعنی !!a{derivation}~$!A_1$، \dots،
  $!A_k = !A$، جیہدے ہر $!A_i$ یا تاں مسلّمہ اے، یا~$\Gamma$
  دا !!a{element} اے، یا استنباط دے قاعدے نال درست اے۔ ہُن
  ایہ سلسلہ ویکھو:
  \[
  !A_1, \dots, !A_k = !A, !B_1, \dots, !B_l = !B.
  \]
% PNBAXD-001: آخری وضاحت وچ B_i دے چھڈے ہوے فارمولا نشان نوں بحال کیتا۔
% PNBAXD-002: تین ممکن جوازاں لئی ماخذ دے صرف قاعدہ آکھن نوں جواز کیتا۔
  ایہ~$!B$ دا $\Gamma \cup \Delta$ توں درست !!{derivation} اے،
  کیوں جے ہر $!B_i = !A$ ہُن اوہو جواز رکھدا اے جیہڑا~$!A_k = !A$
  نوں درست ٹھہراندا اے۔
\end{proof}

دھیان دیو کہ ایس دا خاص مطلب ایہ اے: جے $\Gamma \Proves !A$ تے $!A
\Proves !B$، تاں $\Gamma \Proves !B$۔ ایہ وی نکلدا اے کہ جے $!A_1,
\dots, !A_n \Proves !B$ تے $\Gamma \Proves !A_i$ ہر~$i$ لئی، تاں
$\Gamma \Proves !B$۔

\begin{prop}
\ollabel{prop:incons}
$\Gamma$ اُدوں تے صرف اُدوں متضاد اے جدوں $\Gamma \Proves !A$ ہر~$!A$ لئی۔
\end{prop}

\begin{proof}
مشق۔
\end{proof}

\tagprob{FOL}
\begin{prob}
\olref[fol][axd][ptn]{prop:incons} ثابت کرو۔
\end{prob}
\tagendprob

\tagprob{notFOL}
\begin{prob}
\olref[pl][axd][ptn]{prop:incons} ثابت کرو۔
\end{prob}
\tagendprob

\begin{prop}[محدود گواہی (Compactness)]
\ollabel{prop:proves-compact}
  \begin{enumerate}
  \item جے $\Gamma \Proves !A$، تاں اک محدود ذیلی سیٹ $\Gamma_0
    \subseteq \Gamma$ ایہو جہا اے کہ $\Gamma_0 \Proves !A$۔
  \item جے~$\Gamma$ دا ہر محدود ذیلی سیٹ بے تضاد اے، تاں
    $\Gamma$ بے تضاد اے۔
  \end{enumerate}
\end{prop}

\begin{proof}
  \begin{enumerate}
% PNBAXD-012: پہلی درجے دے اخذ وچ مسوّری قاعدیاں تے اوہناں دی غیر حاضری شرط نوں وی شامل کیتا۔
    \item جے $\Gamma \Proves !A$، تاں !!{formula}s دا اک محدود
      سلسلہ $!A_1$، \dots،~$!A_n$ موجود اے جیہدے لئی
      $!A \ident !A_n$ اے؛ ہر $!A_i$ یا تاں منطقی مسلّمہ اے، یا
      ~$\Gamma$ دا !!a{element} اے، یا پچھلے !!{formula}s توں
      قاعدۂ اثباتِ مقدم نال، یا پہلی درجے دی منطق وچ مسوّراں دے
      مقررہ استنباطی قاعدے نال، نکلدا اے۔ $\Gamma_0$ اوہناں $!A_i$
      دا سیٹ لو جیہڑے~$\Gamma$ وچ نیں۔ تاں اوہ !!{derivation}
      ایسے طرح~$\Gamma_0$ توں وی !!a{derivation} اے: مسوّری قاعدے
      دا خاص مستقل اصل مفروضیاں وچ غیر حاضر سی، ایس لئی چُنے ہوے
      ذیلی سیٹ وچ وی غیر حاضر رہندا اے؛قاعدے دیاں باقی غیر حاضری
      شرطاں وی نہیں بدلّدیاں۔ تے ایس لئی
      $\Gamma_0 \Proves !A$۔
    \item ایہہ (1) دا برعکس بالمخالف اے، اوس خاص حالت وچ جدوں $!A
      \ident \lfalse$۔
\end{enumerate}
\end{proof}

\end{document}
